%------------------------------------------------------------------------------%
\documentclass[fleqn,utf8,aspectratio=169,12pt]{beamer}

\usepackage{integracao_e_series}
\usepackage{cancel}

\title{Séries Telescópicas}

%------------------------------------------------------------------------------%
\begin{document}

\addtitle

%------------------------------------------------------------------------------%
\section{Séries Telescópicas}

%------------------------------------------------------------------------------%
\begin{frame}{Séries Telescópicas}
	Raramente conseguimos fórmulas exatas para a soma de séries

	\pause
	\vspace{\baselineskip}
	\emph{Séries Telescópicas} são um caso especial onde quase todos os termos se anulam
\end{frame}

%------------------------------------------------------------------------------%
\section{Exemplos}

%------------------------------------------------------------------------------%
\begin{question}
	Calcular a soma da série
	\[\sum_{n=1}^{\infty} \left( \frac{1}{n} - \frac{1}{n+1} \right) \]
\end{question}

%------------------------------------------------------------------------------%
\begin{resolution}{Somas Parciais}

  {\setlength{\mathindent}{0cm}
	\begin{align*}
		\onslide<+->{S_n & = \sum_{k=1}^{n} \left( \frac{1}{k} - \frac{1}{k+1} \right)
			\\[3mm]}
		\onslide<+->{& = \left( \frac{1}{1} - \frac{1}{1+1} \right)
			+ \left( \frac{1}{2} - \frac{1}{2+1} \right)
			+ \cdots
			+ \left( \frac{1}{n-1} - \frac{1}{n} \right)
			+ \left( \frac{1}{n} - \frac{1}{n+1} \right)
			\\[3mm]}
		\onslide<+->{&
      = \frac{1}{1} - {\color<4->{blue}\frac{1}{2}}
			+ {\color<4->{blue}\frac{1}{2}} - {\color<5->{red}\frac{1}{3}}
      + {\color<5->{red}\frac{1}{3}}  - {\color<6->{gray}\frac{1}{4}}
			+ \cdots
      + {\color<6->{gray}\frac{1}{n-2}}  - {\color<7->{green}\frac{1}{n-1}}
			+ {\color<7->{green}\frac{1}{n-1}} - {\color<8->{magenta}\frac{1}{n}}
			+ {\color<8->{magenta}\frac{1}{n}} - \frac{1}{n+1}
			\\[3mm]}
		\onslide<9->{& = 1 - \frac{1}{n+1}}
	\end{align*}}
\end{resolution}

%------------------------------------------------------------------------------%
\begin{resolution}{Calculando o Limite}
	\begin{align*}
		\onslide<+->{\sum_{n=1}^{\infty} \left( \frac{1}{n} - \frac{1}{n+1} \right)
                  & = \lim_{n\to\infty} S_n \\[3mm]}
		\onslide<+->{ & = \lim_{n\to\infty} \left( 1 - \frac{1}{n+1} \right)  \\[3mm]}
		\onslide<+->{ & = 1}
	\end{align*}
\end{resolution}

%------------------------------------------------------------------------------%

%------------------------------------------------------------------------------%
\begin{question}

 Calcular a soma da série
	\[
    \sum_{n=1}^{\infty} \frac{1}{\;n(n+1)\;}
  \]

\end{question}

%------------------------------------------------------------------------------%
\begin{resolution}{Pode Não Ser Óbvio}

  Usando frações parciais
  \[
    \frac{1}{n(n+1)} = \frac{1}{n} - \frac{1}{n+1}
  \]

  \pause
  Temos
  \[
    \sum_{n=1}^{\infty} \frac{1}{\;n(n+1)\;}
    \;=\; \sum_{n=1}^{\infty} \left( \frac{1}{n} - \frac{1}{n+1} \right)
    \;=\; 1
  \]
\end{resolution}

%------------------------------------------------------------------------------%


%------------------------------------------------------------------------------%
\section{Lista Mínima}

%------------------------------------------------------------------------------%
\begin{frame}{Lista Mínima}
  Estudar as Seção 6.1 da Apostila

  \vspace{\baselineskip}
  Exercícios: 8a-c, 10, 11a-b

  \vfill
  Atenção: A prova é baseada no livro, não nas apresentações
\end{frame}

%------------------------------------------------------------------------------%
\end{document}
%------------------------------------------------------------------------------%
